% =============================================================================
% teams/science/design.tex  -  Science Team design contribution
%
% Filled in by: Science Team
% This file is included in sections/preliminary_design_description.tex
% under Section 3.2 "Rover Design Description" > "Science Module".
% =============================================================================

% ---- Overview ---------------------------------------------------------------
\paragraph{Overview.}
[Describe the purpose of the science module: what scientific tasks it must
perform during ERC (soil/rock sample analysis, astrobiology, geology assessment,
etc.) and the high-level design concept chosen.]

% ---- Instrument Suite -------------------------------------------------------
\paragraph{Instrument Suite.}
[List and briefly describe each scientific instrument included in the module.
State whether each is COTS or custom-built, and its measurement principle.]

\begin{table}[H]
    \centering
    \caption{Science instrument suite}
    \label{tab:science_instruments}
    \begin{tabular}{L{3cm} L{2cm} L{4cm} L{3cm}}
        \toprule
        \textbf{Instrument}       & \textbf{Type} & \textbf{Measured parameter} & \textbf{Notes} \\
        \midrule
        {[Instrument 1]} & COTS / Custom & [pH / conductivity / ...] & \\
        {[Instrument 2]} & COTS / Custom & [Spectroscopy / ...] & \\
        {[Instrument 3]} & COTS / Custom & [Soil moisture / ...] & \\
        % Add all instruments
        \bottomrule
    \end{tabular}
\end{table}

% ---- Sample Collection System -----------------------------------------------
\paragraph{Sample Collection System.}
[Describe how samples are acquired (drilling, scooping, suction, etc.),
stored, and presented to the instruments. Include a figure if available.]

\begin{figure}[H]
    \centering
    \fbox{\parbox{0.8\textwidth}{\centering\vspace{2cm}
        [Insert science module CAD / photo here]\vspace{2cm}}}
    \caption{Science module -- overview}
    \label{fig:science_module}
\end{figure}

% ---- Data Acquisition & Processing ------------------------------------------
\paragraph{Data Acquisition \& Processing.}
[Describe how measurement data is acquired, processed on-board, and transmitted
to the ground station. Mention calibration procedures.]

% ---- Scientific Methodology -------------------------------------------------
\paragraph{Scientific Methodology.}
[Explain the scientific rationale: what signs of life / habitability / geological
interest the module looks for and why the selected instruments are the best
tools for those measurements in an ERC context.]

% ---- Key Specifications Table -----------------------------------------------
\begin{table}[H]
    \centering
    \caption{Science module key specifications}
    \label{tab:science_specs}
    \begin{tabular}{L{5cm} L{7cm}}
        \toprule
        \textbf{Parameter}         & \textbf{Value / Description} \\
        \midrule
        Number of instruments      & [N] \\
        Total module mass          & [X.X kg] \\
        Power consumption (avg.)   & [X.X W] \\
        Sample volume handled      & [X mL / g] \\
        Data interface             & [UART / SPI / I\textsuperscript{2}C / USB] \\
        \bottomrule
    \end{tabular}
\end{table}
